% GENERATED FILE. Edit canonical lessons, then run npm run generate:books.
% canonical-source-hash: fnv1a64:8b7eb8459e5473e8
% canonical-lessons: PT-C399-rede, PT-C399-navegar, PT-C399-pesquisar, PT-C399-motor-de-busca, PT-C399-hiperligacao

\chapter{Na rede}
\label{ch:pt-na-rede}
\hlchaptermodality{399}

\begin{chapteropening}
\textbf{By the end of this chapter:} \emph{I can get about online: browse, search, use a search engine and follow a link.}

Complete the last lesson of chapter 399: I can get about online: browse, search, use a search engine and follow a link.
\end{chapteropening}

\section[a rede]{a rede — the network; the net}

\label{lesson:PT-C399-rede}



\begin{quote}
Before the new one: say the Portuguese for the attachment, then the Portuguese for to attach.
\end{quote}

\subsection*{You'll want to know: a rede}
\textbf{a rede} — \textquotedblleft{}the network; the net\textquotedblright{}.

\begin{quote}
Não há rede aqui na aldeia.
\end{quote}

\begin{quote}
\emph{There is no signal here in the village.}
\end{quote}

\emph{Estar na rede} — to be online; \emph{as redes sociais} — social media. Also a fishing net.

\begin{cousinweb}
\emph{rede} is inherited from Latin \emph{rēte} \textquotedblleft{}a net\textquotedblright{} (EtymDB, Lewis \& Short), through Old Portuguese \emph{rede}.
\end{cousinweb}

\begin{grammarlens}[title={What you've built}]
One more word for telling, explaining and describing at B1.
\end{grammarlens}

\subsection*{Your turn}
\begin{itemize}
\raggedright
  \item \emph{Say it:} \emph{a rede}
  \item \emph{Say it:} \emph{a rede}, once more
  \item \emph{Say it:} a sentence of your own with \emph{a rede}
  \item \emph{Recall:} say the Portuguese for the attachment, then the Portuguese for to attach, then say \emph{a rede} again
\end{itemize}

\begin{tcolorbox}[breakable,colback=teal!4,colframe=teal!35!black,title={Before you move on}]
What does a rede mean? (\textquotedblleft{}The network; the net\textquotedblright{}.) Say it once more.
\end{tcolorbox}
\section[navegar]{navegar — to browse, to surf; to sail}

\label{lesson:PT-C399-navegar}



\begin{quote}
Before the new one: say the Portuguese for to attach, then the Portuguese for the network.
\end{quote}

\subsection*{You'll want to know: navegar}
\textbf{navegar} — \textquotedblleft{}to browse, to surf; to sail\textquotedblright{}.

\begin{quote}
Passo muito tempo a navegar no telemóvel.
\end{quote}

\begin{quote}
\emph{I spend a lot of time browsing on my phone.}
\end{quote}

First to sail a ship.

\begin{cousinweb}
\emph{navegar} comes from Latin \emph{nāvigāre} \textquotedblleft{}to sail\textquotedblright{} (EtymDB), built on \emph{nāvis} \textquotedblleft{}ship\textquotedblright{} (Lewis \& Short).
\end{cousinweb}

\begin{grammarlens}[title={What you've built}]
One more word for telling, explaining and describing at B1.
\end{grammarlens}

\subsection*{Your turn}
\begin{itemize}
\raggedright
  \item \emph{Say it:} \emph{navegar}
  \item \emph{Say it:} \emph{navegar}, once more
  \item \emph{Say it:} a sentence of your own with \emph{navegar}
  \item \emph{Recall:} say the Portuguese for to attach, then the Portuguese for the network, then say \emph{navegar} again
\end{itemize}

\begin{tcolorbox}[breakable,colback=teal!4,colframe=teal!35!black,title={Before you move on}]
What does navegar mean? (\textquotedblleft{}To browse, to surf; to sail\textquotedblright{}.) Say it once more.
\end{tcolorbox}
\section[pesquisar]{pesquisar — to search, to look up; to research}

\label{lesson:PT-C399-pesquisar}



\begin{quote}
Before the new one: say the Portuguese for the network, then the Portuguese for to browse.
\end{quote}

\subsection*{You'll want to know: pesquisar}
\textbf{pesquisar} — \textquotedblleft{}to search, to look up; to research\textquotedblright{}.

\begin{quote}
Pesquisei o horário do comboio no telemóvel.
\end{quote}

\begin{quote}
\emph{I looked up the train times on my phone.}
\end{quote}

\emph{Fazer uma pesquisa} — to do a search, to do research.

\begin{cousinweb}
\emph{pesquisar} is built on \emph{pesquisa} \textquotedblleft{}search\textquotedblright{}, which EtymDB traces through Spanish \emph{pesquisa} to Latin \emph{perquīrere} \textquotedblleft{}to make diligent search for\textquotedblright{} (EtymDB, Lewis \& Short), built on \emph{quaerere} \textquotedblleft{}to seek\textquotedblright{}.
\end{cousinweb}

\begin{grammarlens}[title={What you've built}]
One more word for telling, explaining and describing at B1.
\end{grammarlens}

\subsection*{Your turn}
\begin{itemize}
\raggedright
  \item \emph{Say it:} \emph{pesquisar}
  \item \emph{Say it:} \emph{pesquisar}, once more
  \item \emph{Say it:} a sentence of your own with \emph{pesquisar}
  \item \emph{Recall:} say the Portuguese for the network, then the Portuguese for to browse, then say \emph{pesquisar} again
\end{itemize}

\begin{tcolorbox}[breakable,colback=teal!4,colframe=teal!35!black,title={Before you move on}]
What does pesquisar mean? (\textquotedblleft{}To search, to look up; to research\textquotedblright{}.) Say it once more.
\end{tcolorbox}
\section[o motor de busca]{o motor de busca — the search engine}

\label{lesson:PT-C399-motor-de-busca}



\begin{quote}
Before the new one: say the Portuguese for to browse, then the Portuguese for to search.
\end{quote}

\subsection*{You'll want to know: o motor de busca}
\textbf{o motor de busca} — \textquotedblleft{}the search engine\textquotedblright{}.

\begin{quote}
Escrevi o nome do restaurante no motor de busca.
\end{quote}

\begin{quote}
\emph{I typed the name of the restaurant into the search engine.}
\end{quote}

Also \emph{o motor de pesquisa}.

\begin{cousinweb}
\emph{motor} comes through French \emph{moteur} from Latin \emph{mōtor} \textquotedblleft{}a mover\textquotedblright{} (EtymDB), on \emph{movēre} \textquotedblleft{}to move\textquotedblright{} (Lewis \& Short); \emph{busca} \textquotedblleft{}search\textquotedblright{} is built on \emph{buscar} \textquotedblleft{}to look for\textquotedblright{}, whose origin is not settled.
\end{cousinweb}

\begin{grammarlens}[title={What you've built}]
One more word for telling, explaining and describing at B1.
\end{grammarlens}

\subsection*{Your turn}
\begin{itemize}
\raggedright
  \item \emph{Say it:} \emph{o motor de busca}
  \item \emph{Say it:} \emph{o motor de busca}, once more
  \item \emph{Say it:} a sentence of your own with \emph{o motor de busca}
  \item \emph{Recall:} say the Portuguese for to browse, then the Portuguese for to search, then say \emph{o motor de busca} again
\end{itemize}

\begin{tcolorbox}[breakable,colback=teal!4,colframe=teal!35!black,title={Before you move on}]
What does o motor de busca mean? (\textquotedblleft{}The search engine\textquotedblright{}.) Say it once more.
\end{tcolorbox}
\section[a hiperligação]{a hiperligação — the web link, the hyperlink}

\label{lesson:PT-C399-hiperligacao}



\begin{quote}
Before the new one: say the Portuguese for to search, then the Portuguese for the search engine.
\end{quote}

\subsection*{You'll want to know: a hiperligação}
\textbf{a hiperligação} — \textquotedblleft{}the web link, the hyperlink\textquotedblright{}.

\begin{quote}
Cliquei na hiperligação e abriu uma página nova.
\end{quote}

\begin{quote}
\emph{I clicked on the link and a new page opened.}
\end{quote}

European Portuguese, formal; in speech most people say \emph{o link}.

\begin{cousinweb}
\emph{hiperligação} is built from Greek \emph{hypér} \textquotedblleft{}over, beyond\textquotedblright{} and \emph{ligação} \textquotedblleft{}connection\textquotedblright{}, on \emph{ligar}, inherited from Latin \emph{ligāre} \textquotedblleft{}to tie\textquotedblright{} (Lewis \& Short); it translates English \emph{hyperlink}.
\end{cousinweb}

\begin{grammarlens}[title={What you've built}]
One more word for telling, explaining and describing at B1.
\end{grammarlens}

\subsection*{Your turn}
\begin{itemize}
\raggedright
  \item \emph{Say it:} \emph{a hiperligação}
  \item \emph{Say it:} \emph{a hiperligação}, once more
  \item \emph{Say it:} a sentence of your own with \emph{a hiperligação}
  \item \emph{Recall:} say the Portuguese for to search, then the Portuguese for the search engine, then say \emph{a hiperligação} again
\end{itemize}

\begin{tcolorbox}[breakable,colback=teal!4,colframe=teal!35!black,title={Before you move on}]
What does a hiperligação mean? (\textquotedblleft{}The web link, the hyperlink\textquotedblright{}.) Say it once more.
\end{tcolorbox}
